% Autor szablonu: dr Roman Czapla
% Kompilacja: pdfLaTeX (dwukrotnie, aby zaktualizować spis treści).
% Wymagane logo uczelni: images/logo.png (ścieżka względem pliku głównego).
\documentclass[11pt,a4paper,oneside]{article}
\usepackage[utf8]{inputenc}
\usepackage[T1]{fontenc}
\usepackage[polish,provide=*]{babel}
\usepackage{lmodern,microtype}
\usepackage[margin=25mm,headheight=14pt]{geometry}
\usepackage{xcolor,graphicx,tabularx,booktabs,array,enumitem,fancyhdr,titlesec}
\usepackage[hidelinks]{hyperref}
\definecolor{accent}{HTML}{344B59}
\definecolor{hintcolor}{HTML}{58636A}
\hypersetup{pdfauthor={},pdftitle={Projekt inżynierski - Dokumentacja użytkowa}}
\setlength{\parindent}{0pt}
\setlength{\parskip}{6pt}
\renewcommand{\arraystretch}{1.3}
\setlength{\emergencystretch}{2em}
\setlist[itemize]{leftmargin=*,itemsep=2pt,topsep=3pt}
\titleformat{\section}{\Large\bfseries\color{accent}}{\thesection.}{0.5em}{}
\titleformat{\subsection}{\large\bfseries\color{accent}}{\thesubsection.}{0.5em}{}
\pagestyle{fancy}
\fancyhf{}
\fancyhead[L]{\small Projekt inżynierski}
\fancyhead[R]{\small Dokumentacja użytkowa}
\fancyfoot[C]{\small\thepage}
\renewcommand{\headrulewidth}{0.3pt}
\newcolumntype{Y}{>{\raggedright\arraybackslash}X}

% Podpowiedzi są widoczne podczas wypełniania szablonu.
% Przed oddaniem: uzupełnij pola, zastąp przykłady własną treścią
% i zmień poniższy przełącznik na \podpowiedzifalse.
\newif\ifpodpowiedzi
\podpowiedzitrue
%\podpowiedzifalse
\newcommand{\podpowiedz}[1]{\ifpodpowiedzi\par\begingroup\small\color{hintcolor}\textit{Wskazówka: #1}\par\endgroup\fi}
\newcommand{\pole}[1]{\textnormal{[#1]}}

% Dane dokumentu - uzupełnij w jednym miejscu.
% Ścieżkę logo można zmienić poniżej.
\newcommand{\pliklogo}{images/logo.png}
\newcommand{\logouczelni}{%
  \includegraphics[width=0.8\textwidth,height=40mm,keepaspectratio]{\pliklogo}%
}
\newcommand{\tytulprojektu}{\pole{Tytuł projektu}}
\newcommand{\rokakademicki}{\pole{2026/2027}}
\newcommand{\semestrprojektu}{\pole{zimowy / letni}}
\newcommand{\datadokumentu}{\pole{dd.mm.rrrr}}
\newcommand{\wersjadokumentu}{\pole{1.0}}
% Zespół dwuosobowy tylko za zgodą dyrekcji: zmień na \trzeciaosobafalse.
\newif\iftrzeciaosoba
\trzeciaosobatrue
\newcommand{\studentA}{\pole{Imię i nazwisko osoby 1}}
\newcommand{\studentB}{\pole{Imię i nazwisko osoby 2}}
\newcommand{\studentC}{\pole{Imię i nazwisko osoby 3}}
\newcommand{\albumA}{\pole{numer albumu}}
\newcommand{\albumB}{\pole{numer albumu}}
\newcommand{\albumC}{\pole{numer albumu}}
\newcommand{\prowadzacyA}{\pole{Stopień lub tytuł, imię i nazwisko}}
\newcommand{\prowadzacyB}{\pole{Stopień lub tytuł, imię i nazwisko}}

\begin{document}
\begin{titlepage}
\centering
{\large Uniwersytet Komisji Edukacji Narodowej w Krakowie\par}
{\normalsize Instytut Bezpieczeństwa i Informatyki\par}
\vspace{8mm}
\logouczelni\par
\vspace{10mm}
{\Large\bfseries\color{accent} PROJEKT INŻYNIERSKI\par}
\vspace{3mm}
{\large Dokumentacja użytkowa\par}
\vspace{15mm}
{\LARGE\bfseries\tytulprojektu\par}
\vspace{15mm}
{\small ZESPÓŁ PROJEKTOWY\par}
\vspace{3mm}
\begin{tabular}{ll}
\studentA & Nr albumu: \albumA\\
\studentB & Nr albumu: \albumB
\iftrzeciaosoba\\
\studentC & Nr albumu: \albumC\fi
\end{tabular}
\vspace{12mm}

{\small PROWADZĄCY ZAJĘCIA PROJEKTOWE\par}
\vspace{3mm}
\prowadzacyA\par
\prowadzacyB\par
\vfill
Rok akademicki: \rokakademicki\quad Semestr: \semestrprojektu\par
Data dokumentu: \datadokumentu\quad Wersja: \wersjadokumentu\par
\vspace{6mm}
Kraków\par

\end{titlepage}

\tableofcontents
\podpowiedz{Dokument jest skierowany do osoby korzystającej z rezultatu projektu. Zakres uzgodnij z prowadzącymi i dostosuj do odbiorcy: użytkownika aplikacji, operatora prototypu, odbiorcy animacji albo programisty korzystającego z biblioteki. Usuń podsekcje nieadekwatne. Nie powielaj opisu architektury ani instrukcji instalacji, jeśli znajdują się w osobnych dokumentach. Uzupełnij pola, usuń przykłady i wyłącz podpowiedzi przed oddaniem.}

\section{Szczegółowa dokumentacja użytkowa}
\subsection{Przeznaczenie i odbiorcy}
\podpowiedz{Wyjaśnij, do czego służy rozwiązanie, kto może z niego korzystać i jakie zadania umożliwia. Podaj wersję projektu, której dotyczy instrukcja. Używaj języka zrozumiałego dla wskazanego odbiorcy.}
\textbf{Przeznaczenie:} \pole{opis.}\par
\textbf{Odbiorcy:} \pole{grupa użytkowników i potrzebne umiejętności.}\par
\textbf{Wersja projektu:} \pole{oznaczenie i data.}


\subsection{Rozpoczęcie korzystania}
\podpowiedz{Podaj wymagania istotne dla użytkownika i sposób rozpoczęcia pracy: adres aplikacji, uruchamiany plik, odtwarzacz, urządzenie lub punkt wejścia do biblioteki. Wskaż potrzebne konto, uprawnienia i dane wejściowe, jeśli występują. Odwołaj się do dokumentacji instalacyjnej, jeśli jest wymagana; nie opisuj środowiska deweloperskiego.}
\pole{Wymagania i sposób rozpoczęcia korzystania.}

\subsection{Opis funkcji i sposobu obsługi}
\podpowiedz{Opisz dostępne funkcje oraz sposób ich użycia, grupując je według zadań użytkownika. Wyjaśnij istotne elementy interfejsu, polecenia i ustawienia. Dla animacji może to być sposób odtwarzania i korzystania z przekazanych formatów; dla prototypu - sterowanie; dla biblioteki - publiczny interfejs i przykłady wywołań.}
\pole{Opis funkcji i czynności użytkownika.}
\podpowiedz{Dodaj czytelne zrzuty ekranu lub ilustracje tam, gdzie pomagają wykonać zadanie. Każdy rysunek powinien mieć podpis i odwołanie w tekście przez \texttt{\textbackslash ref}. Nie zastępuj instrukcji samymi obrazami.}
% Przykład - odkomentuj po dodaniu pliku:
% \begin{figure}[htbp]
%   \centering
%   \includegraphics[width=0.9\linewidth]{images/ekran.png}
%   \caption{Widok funkcji z perspektywy użytkownika.}
%   \label{fig:ekran}
% \end{figure}

\subsection{Przykładowe działanie}
\podpowiedz{Przedstaw konkretne scenariusze najważniejszych zadań. Dla każdego podaj cel, warunki początkowe, przykładowe dane, kolejne kroki i oczekiwany rezultat. Przykład powinien być możliwy do odtworzenia. Powielaj poniższy blok według potrzeb.}
\subsubsection{\pole{Nazwa scenariusza}}
\textbf{Cel:} \pole{zadanie użytkownika.}\par
\textbf{Warunki i dane:} \pole{stan początkowy i przykładowe dane.}
\begin{enumerate}
\item \pole{czynność użytkownika i reakcja rozwiązania.}
\item \pole{kolejna czynność i reakcja.}
\item \pole{zakończenie zadania.}
\end{enumerate}
\textbf{Oczekiwany rezultat:} \pole{efekt i sposób jego rozpoznania.}

\subsection{Dane wejściowe i wyniki}
\podpowiedz{Jeśli dotyczy: opisz obsługiwane formaty, wymagane pola, dopuszczalne wartości i ograniczenia danych. Wyjaśnij zapis, eksport, odtwarzanie lub interpretację wyników. Wskaż czynności mogące nadpisać lub usunąć dane. Dla animacji możesz opisać formaty plików wynikowych i parametry odtwarzania.}
\pole{Przygotowanie danych i korzystanie z wyników.}

\subsection{Problemy i komunikaty}
\podpowiedz{Opisz typowe problemy z perspektywy użytkownika: jak je rozpoznać, co oznaczają i jakie działania podjąć. Podaj rzeczywiste komunikaty, jeśli występują. Dodaj wiersze według potrzeb.}
\begin{tabularx}{\textwidth}{@{}Y Y Y@{}}
\toprule
Objaw / komunikat & Możliwa przyczyna & Działanie użytkownika\\
\midrule
\pole{objaw} & \pole{przyczyna} & \pole{rozwiązanie}\\
\pole{objaw} & \pole{przyczyna} & \pole{rozwiązanie}\\
\bottomrule
\end{tabularx}

\subsection{Ograniczenia i zalecenia}
\podpowiedz{Wskaż ograniczenia mające wpływ na korzystanie z projektu oraz praktyczne zalecenia. Jeśli użytkownik korzysta z funkcji AI, wyjaśnij charakter wyników, ich ograniczenia i potrzebę weryfikacji tam, gdzie ma to znaczenie. Nie dodawaj informacji o AI, jeśli nie jest częścią rozwiązania i nie wpływa na użytkowanie.}
\pole{Ograniczenia i zalecenia dla użytkownika.}

\subsection{Materiały pomocnicze}
\podpowiedz{Wskaż przydatne instrukcje, przykładowe pliki oraz obowiązkowe nagranie pokazujące działanie lub rezultat oddanej wersji projektu. Podaj nazwy i lokalizację materiałów. Źródła zewnętrzne przywołane w instrukcji opisz tytułem, adresem i datą dostępu. Jeśli materiały zostały już wymienione w sprawozdaniu końcowym, możesz odwołać się do odpowiedniej sekcji. Podsekcję można usunąć, jeśli nie wnosi dodatkowych informacji; nie znosi to obowiązku przekazania nagrania i wskazania go w wykazie załączników sprawozdania końcowego.}
\pole{Wykaz materiałów i ich lokalizacja lub odwołanie do sprawozdania końcowego.}
\end{document}
